\documentclass[UTF8,a4paper,11pt]{ctexart}

\usepackage[margin=2.3cm]{geometry}
\usepackage{xcolor}
\usepackage{graphicx}
\usepackage{enumitem}
\usepackage{array}
\usepackage{longtable}
\usepackage{booktabs}
\usepackage{titlesec}
\usepackage{fancyhdr}
\usepackage{tcolorbox}
\usepackage{hyperref}
\usepackage{listings}

\tcbuselibrary{skins,breakable}

\definecolor{accentblue}{HTML}{1F5FAD}
\definecolor{accentcyan}{HTML}{0A7D96}
\definecolor{warnorange}{HTML}{B86E00}
\definecolor{dangerred}{HTML}{B42020}
\definecolor{successgreen}{HTML}{1F8A5A}
\definecolor{softgray}{HTML}{F4F6FA}
\definecolor{bordergray}{HTML}{D7DCE5}
\definecolor{codebg}{HTML}{F7F8FC}

\hypersetup{
  colorlinks=true,
  linkcolor=accentblue,
  urlcolor=accentcyan,
  citecolor=accentblue,
  pdftitle={废气综合管理平台 Docker 部署手册},
  pdfauthor={智洁园区项目组}
}

\titleformat{\section}
  {\Large\bfseries\color{accentblue}}
  {\thesection}{0.6em}{}
  [\vspace{2pt}{\color{bordergray}\titlerule[0.8pt]}]
\titleformat{\subsection}
  {\large\bfseries\color{accentcyan}}
  {\thesubsection}{0.5em}{}
\titleformat{\subsubsection}
  {\normalsize\bfseries}
  {\thesubsubsection}{0.5em}{}

\pagestyle{fancy}
\fancyhf{}
\fancyhead[L]{\small\color{accentblue}\textbf{废气综合管理平台}}
\fancyhead[R]{\small\color{accentcyan}Docker 部署手册 v2.0}
\fancyfoot[C]{\small\thepage}
\renewcommand{\headrulewidth}{0.4pt}
\renewcommand{\footrulewidth}{0pt}

\lstdefinestyle{shellstyle}{
  basicstyle=\ttfamily\small,
  backgroundcolor=\color{codebg},
  frame=single,
  rulecolor=\color{bordergray},
  framesep=4pt,
  xleftmargin=6pt,
  xrightmargin=4pt,
  breaklines=true,
  breakatwhitespace=true,
  showstringspaces=false,
  keywordstyle=\color{accentblue}\bfseries,
  commentstyle=\color{gray}\itshape,
  stringstyle=\color{warnorange},
  aboveskip=6pt,
  belowskip=6pt,
}
\lstset{style=shellstyle}
\lstdefinelanguage{dockerfile}{
  keywords={FROM,RUN,CMD,ENV,COPY,ADD,WORKDIR,EXPOSE,VOLUME,ENTRYPOINT,USER,ARG,LABEL,HEALTHCHECK,AS},
  sensitive=true,
  morecomment=[l]{\#},
  morestring=[b]"
}

\newtcolorbox{infobox}[1][]{
  colback=softgray, colframe=accentblue, coltitle=white,
  fonttitle=\bfseries, title=#1,
  sharp corners=downhill, boxrule=0.6pt, left=8pt, right=8pt, top=6pt, bottom=6pt,
  breakable
}
\newtcolorbox{successbox}[1][]{
  colback=softgray, colframe=successgreen, coltitle=white,
  fonttitle=\bfseries, title=#1,
  sharp corners=downhill, boxrule=0.6pt, left=8pt, right=8pt, top=6pt, bottom=6pt,
  breakable
}
\newtcolorbox{warnbox}[1][]{
  colback=softgray, colframe=warnorange, coltitle=white,
  fonttitle=\bfseries, title=#1,
  sharp corners=downhill, boxrule=0.6pt, left=8pt, right=8pt, top=6pt, bottom=6pt,
  breakable
}
\newtcolorbox{dangerbox}[1][]{
  colback=softgray, colframe=dangerred, coltitle=white,
  fonttitle=\bfseries, title=#1,
  sharp corners=downhill, boxrule=0.6pt, left=8pt, right=8pt, top=6pt, bottom=6pt,
  breakable
}
\newtcolorbox{stepbox}[1][]{
  colback=white, colframe=accentblue, coltitle=white,
  fonttitle=\bfseries\large, title=#1,
  sharp corners=downhill, boxrule=1pt, left=10pt, right=10pt, top=8pt, bottom=8pt,
  breakable
}

\renewcommand{\arraystretch}{1.25}
\setlength{\tabcolsep}{6pt}
\newcolumntype{L}[1]{>{\raggedright\arraybackslash}p{#1}}

\begin{document}

% --- 封面 ---
\begin{titlepage}
\centering
\vspace*{2.5cm}
{\color{accentblue}\rule{14cm}{1.2pt}}\par
\vspace{0.4cm}
{\Huge\bfseries 智洁园区\,·\,废气综合管理平台\par}
\vspace{0.6cm}
{\LARGE\bfseries\color{accentblue} 系统部署手册\par}
\vspace{0.2cm}
{\Large\color{accentcyan}（Docker 容器化版本）\par}
\vspace{0.3cm}
{\color{accentblue}\rule{14cm}{1.2pt}}\par
\vspace{1.6cm}
\begin{tcolorbox}[colback=softgray, colframe=accentblue, sharp corners, boxrule=0.6pt, width=12cm]
\centering
\large\textbf{4 步上线} \quad $|$ \quad
\textbf{4 个容器} \quad $|$ \quad
\textbf{1 个入口}\par
\vspace{4pt}
\small 全栈 Docker 化\quad·\quad支持离线交付\quad·\quad支持云端部署
\end{tcolorbox}
\vspace{1.6cm}
\begin{tabular}{rl}
  \textbf{文档版本} & v2.0 (Docker 版) \\
  \textbf{发布日期} & 2026-04-22 \\
  \textbf{适用版本} & admin-frontend 1.0.0 / admin-backend 0.1.0 \\
  \textbf{编排方式} & Docker Compose v2 \\
  \textbf{编写团队} & 智洁园区项目组 \\
\end{tabular}
\vspace{2cm}
\end{titlepage}

\tableofcontents
\newpage

% ========== 1 项目概述 ==========
\section{项目概述}

\subsection{Docker 化方案}

本平台已完成全栈容器化，4 个核心服务全部封装为独立 Docker 镜像，
通过 \texttt{docker compose} 一键编排启动。相比第一版手册（v1.0）依赖
Python/Node/Nginx 多套环境逐一安装的方式，本版本对部署人员的环境要求大幅简化：
\textbf{只要有 Docker 与 Docker Compose，4 步即可完成全部部署}。

\subsection{容器服务一览}

\begin{longtable}{L{3.6cm} L{3.0cm} L{1.6cm} L{6.0cm}}
\toprule
\textbf{服务} & \textbf{容器名} & \textbf{端口} & \textbf{职责} \\
\midrule
\endhead
集成学习预测服务      & wg-ensemble        & \textbf{8000} & VOC 多模型集成预测 (DLinear+XGBoost) \\
VOCs 实时推理服务     & wg-vocs-server     & \textbf{8001} & 传感器接入、Seq2Seq 预测、SSE 推送 \\
管理端后端 (FastAPI)  & wg-admin-backend   & \textbf{8003} & 大屏聚合 API、JWT 鉴权、RAG \\
管理端前端 (Nginx)    & wg-admin-frontend  & \textbf{3001} & Vue3 + Three.js 可视化大屏 \\
\bottomrule
\end{longtable}

\begin{successbox}[最终入口]
4 个容器加入同一个 \texttt{wg-net} 桥接网络，容器内部通过\textbf{服务名互访}
（如 \texttt{http://vocs-server:8001}），上云后无需修改任何配置。
\textbf{用户最终只需访问一个地址：}\texttt{http://<服务器IP>:3001}
\end{successbox}

\subsection{镜像清单}

\begin{longtable}{L{4.5cm} L{2.5cm} L{2.5cm} L{4.0cm}}
\toprule
\textbf{镜像名} & \textbf{基础镜像} & \textbf{预估大小} & \textbf{Dockerfile 位置} \\
\midrule
\endhead
waste-gas/ensemble:latest        & python:3.10-slim    & \textasciitilde{}1.6 GB & \texttt{models/new\_VOC/ensemble\_docker/} \\
waste-gas/vocs-server:latest     & python:3.10-slim    & \textasciitilde{}1.4 GB & \texttt{./Dockerfile} \\
waste-gas/admin-backend:latest   & python:3.11-slim    & \textasciitilde{}2.2 GB & \texttt{admin/backend/Dockerfile} \\
waste-gas/admin-frontend:latest  & nginx:1.27-alpine   & \textasciitilde{}45  MB & \texttt{admin/frontend/Dockerfile} \\
\bottomrule
\end{longtable}

% ========== 2 环境配置要求 ==========
\section{环境配置要求}

\subsection{硬件最低配置}

\begin{longtable}{L{3.0cm} L{4.5cm} L{5.5cm}}
\toprule
\textbf{项目} & \textbf{最低（演示）} & \textbf{推荐（生产）} \\
\midrule
\endhead
CPU      & 4 核 (x86\_64)            & 8 核以上，支持 AVX2 \\
内存      & 8 GB                      & 16 GB 及以上 \\
硬盘      & 20 GB SSD（含镜像缓存）    & 100 GB SSD（含日志/备份） \\
GPU      & 非必须（CPU 推理已可）      & NVIDIA T4/A10（可选加速） \\
网络      & 千兆内网                  & 千兆内网 + 公网备用 \\
\bottomrule
\end{longtable}

\subsection{软件依赖（部署机器）}

\begin{infobox}[只需两样东西]
\textbf{Docker Engine} $\geq$ 24.0  \quad \textbf{+}  \quad
\textbf{Docker Compose} $\geq$ v2.20（内置 \texttt{docker compose} 子命令）

\smallskip
\textbf{不再需要}：手动安装 Python / Node / Nginx / 各种依赖包 / 模型环境。
\end{infobox}

\subsubsection{支持的操作系统}

\begin{itemize}[leftmargin=1.4em,itemsep=2pt]
  \item Linux：Ubuntu 22.04 LTS / Debian 12 / CentOS 9 / Rocky Linux 9（推荐）
  \item Windows：Windows 10/11 + WSL2 + Docker Desktop
  \item macOS：macOS 12+ + Docker Desktop（仅限本地演示）
\end{itemize}

\subsubsection{Docker 安装速查}

\begin{lstlisting}[language=bash]
# Ubuntu / Debian 一键脚本
curl -fsSL https://get.docker.com | bash
sudo usermod -aG docker $USER
newgrp docker

# 验证
docker --version              # 应输出 Docker version 24.x.x
docker compose version        # 应输出 Docker Compose version v2.x.x
\end{lstlisting}

\subsection{端口占用要求}

部署前请确认以下端口在宿主机上空闲（被占用时改 \texttt{docker-compose.prod.yml} 端口映射即可）：

\begin{longtable}{L{2.0cm} L{4.5cm} L{6.5cm}}
\toprule
\textbf{端口} & \textbf{对应服务} & \textbf{是否必须对外暴露} \\
\midrule
\endhead
\textbf{3001} & 管理端前端 (Nginx)         & \textbf{是} —— 用户唯一入口 \\
8003          & 管理端后端 (FastAPI)        & 可选（容器内已通过服务名互访） \\
8001          & VOCs 推理服务              & 可选（用于现场设备直接 POST 数据） \\
8000          & 集成学习预测服务            & 可选（仅供 admin-backend 调用） \\
\bottomrule
\end{longtable}

% ========== 3 Docker 部署步骤（4 步） ==========
\section{Docker 部署步骤（4 步上线）}

\begin{successbox}[流程总览]
\textbf{Step 1}：克隆代码 / 解包交付物 \quad$\rightarrow$\quad
\textbf{Step 2}：复制并修改 \texttt{.env} \quad$\rightarrow$\quad
\textbf{Step 3}：构建镜像 \quad$\rightarrow$\quad
\textbf{Step 4}：启动并验证
\end{successbox}

% ----- Step 1 -----
\begin{stepbox}[Step 1：获取交付物]
\textbf{方式 A：从 Git 克隆（在线）}

\begin{lstlisting}[language=bash]
git clone <your-repo-url> waste-gas
cd waste-gas
\end{lstlisting}

\textbf{方式 B：从离线包解压（无网环境）}

\begin{lstlisting}[language=bash]
# 1) 解压源码包
tar -xzf waste-gas-bundle.tar.gz
cd waste-gas

# 2) 加载预先导出的镜像（约 4.5 GB）
gunzip -c waste-gas-images.tar.gz | docker load
\end{lstlisting}

解压/克隆后，根目录应包含：
\texttt{docker-compose.prod.yml}、\texttt{.env.example}、\texttt{Dockerfile}、
\texttt{vocs\_server.py}、\texttt{models/}、\texttt{admin/}、
\texttt{models/new\_VOC/ensemble\_docker/}。
\end{stepbox}

% ----- Step 2 -----
\begin{stepbox}[Step 2：配置环境变量]
\begin{lstlisting}[language=bash]
cp .env.example .env
# 用任意编辑器打开 .env，按需修改：
#   ADMIN_PASSWORD=改成强密码
#   ADMIN_JWT_SECRET=改成长随机串
#   PG_HOST / PG_USER / PG_PASSWORD（如使用自有数据库）
\end{lstlisting}

\begin{infobox}[最小可用]
若仅作演示，\textbf{可以不修改}直接进入下一步——默认账户 \texttt{admin / admin123456}，
PG 默认指向项目组开发库，所有功能立即可用。
\end{infobox}
\end{stepbox}

% ----- Step 3 -----
\begin{stepbox}[Step 3：一键构建镜像]
\begin{lstlisting}[language=bash]
docker compose -f docker-compose.prod.yml build
\end{lstlisting}

首次构建约 \textbf{8-15 分钟}（主要在拉取 torch、sentence-transformers、
node\_modules）。后续修改代码再 build，多数层会命中缓存，1-3 分钟完成。

\textbf{如使用离线镜像（Step 1 方式 B）}：跳过本步，直接进入 Step 4。
\end{stepbox}

% ----- Step 4 -----
\begin{stepbox}[Step 4：启动并验证]
\begin{lstlisting}[language=bash]
docker compose -f docker-compose.prod.yml up -d
docker compose -f docker-compose.prod.yml ps   # 4 个容器都应是 Up
\end{lstlisting}

\textbf{快速健康检查：}

\begin{lstlisting}[language=bash]
curl http://localhost:8000/health        # ensemble
curl http://localhost:8001/status        # vocs-server
curl http://localhost:8003/api/health    # admin-backend
curl http://localhost:3001/healthz       # admin-frontend
\end{lstlisting}

\textbf{完成！}浏览器访问 \texttt{http://<服务器IP>:3001}，
使用 \texttt{admin / admin123456} 登录即可看到大屏。
\end{stepbox}

\subsection{常用运维命令}

\begin{lstlisting}[language=bash]
# 查看实时日志
docker compose -f docker-compose.prod.yml logs -f admin-backend

# 重启某个服务
docker compose -f docker-compose.prod.yml restart admin-frontend

# 修改前端代码后单独重建
docker compose -f docker-compose.prod.yml build admin-frontend
docker compose -f docker-compose.prod.yml up -d admin-frontend

# 全部停止（保留数据）
docker compose -f docker-compose.prod.yml down

# 全部停止并删除数据卷（慎用）
docker compose -f docker-compose.prod.yml down -v
\end{lstlisting}

\subsection{云端部署 / 离线交付}

\subsubsection{方式 A：导出镜像 tar 包（适合无网环境）}

\begin{lstlisting}[language=bash]
# 在有网机器上构建好后导出
docker save \
  waste-gas/ensemble:latest \
  waste-gas/vocs-server:latest \
  waste-gas/admin-backend:latest \
  waste-gas/admin-frontend:latest \
  | gzip > waste-gas-images.tar.gz

# 把 tar.gz + 项目源码 + .env + models/ + vocs_realtime_data/
# 一起拷贝到目标服务器，按 Step 1 方式 B 加载即可
\end{lstlisting}

\subsubsection{方式 B：推送到镜像仓库（阿里云 ACR / Harbor / Docker Hub）}

\begin{lstlisting}[language=bash]
# 1) 登录仓库
docker login registry.cn-hangzhou.aliyuncs.com

# 2) 重打 tag
docker tag waste-gas/admin-frontend:latest \
  registry.cn-hangzhou.aliyuncs.com/<命名空间>/admin-frontend:1.0.0
# (其余 3 个镜像同理)

# 3) 推送
docker push registry.cn-hangzhou.aliyuncs.com/<命名空间>/admin-frontend:1.0.0
\end{lstlisting}

之后将 \texttt{docker-compose.prod.yml} 中的 \texttt{build:} 块替换为
\texttt{image:}，目标服务器执行 \texttt{docker compose pull \&\& docker compose up -d} 即可。

% ========== 4 依赖包清单 ==========
\section{依赖包清单（已封装在镜像内）}

\begin{infobox}[说明]
以下依赖均已固化在对应镜像中，部署人员\textbf{无需手动安装}。
本节仅供二次开发或镜像审计参考。
\end{infobox}

\subsection{ensemble 镜像（python:3.10-slim）}

\textbf{Dockerfile 位置：}\texttt{models/new\_VOC/ensemble\_docker/Dockerfile}

\begin{longtable}{L{3.5cm} L{2.0cm} L{7.5cm}}
\toprule
\textbf{包名} & \textbf{版本} & \textbf{说明} \\
\midrule
\endhead
fastapi          & >=0.104.1 & Web 框架 \\
uvicorn          & >=0.24.0  & ASGI 服务器 \\
pydantic         & >=2.4.0   & 数据校验 \\
numpy            & >=1.26.0  & 数值计算 \\
pandas           & >=2.0.3   & 时序数据处理 \\
scikit-learn     & >=1.3.0   & 集成学习底层依赖 \\
torch            & >=2.0.0   & DLinear 模型推理 \\
\bottomrule
\end{longtable}

\subsection{vocs-server 镜像（python:3.10-slim）}

\textbf{Dockerfile 位置：}\texttt{./Dockerfile}（项目根目录）

\begin{longtable}{L{3.5cm} L{2.0cm} L{7.5cm}}
\toprule
\textbf{包名} & \textbf{版本} & \textbf{说明} \\
\midrule
\endhead
fastapi          & 0.104.1 & Web 框架 \\
uvicorn          & 0.24.0  & ASGI 服务器 \\
torch            & 2.7.0   & Seq2Seq 预测模型 \\
numpy            & >=1.26.0 & 数值计算 \\
pandas           & >=2.0.3  & 时序数据处理 \\
pydantic         & 2.4.0    & 数据校验 \\
requests         & 2.31.0   & 同步 HTTP 客户端（测试脚本） \\
sseclient        & 0.0.27   & SSE 测试客户端 \\
scikit-learn     & 1.7.0    & StandardScaler 反归一化 \\
\bottomrule
\end{longtable}

\subsection{admin-backend 镜像（python:3.11-slim）}

\textbf{Dockerfile 位置：}\texttt{admin/backend/Dockerfile}

\begin{longtable}{L{3.8cm} L{2.0cm} L{7.2cm}}
\toprule
\textbf{包名} & \textbf{版本} & \textbf{说明} \\
\midrule
\endhead
fastapi                   & 0.116.1 & Web 框架 \\
uvicorn                   & 0.35.0  & ASGI 服务器 \\
httpx                     & 0.28.1  & 异步 HTTP，代理到 vocs/ensemble \\
pyjwt                     & 2.10.1  & 管理员 JWT \\
pydantic                  & 2.11.7  & 数据模型 \\
psycopg2-binary           & 2.9.10  & PostgreSQL 驱动 \\
python-dotenv             & 1.2.2   & .env 加载 \\
sentence-transformers     & 2.2.2   & RAG 向量化 \\
torch (CPU 版)            & 2.3.1   & 向量模型推理 \\
openai                    & 2.31.0  & 大模型接口 \\
numpy                     & 1.26.4  & 向量运算 \\
pandas                    & 3.0.2   & 结构化数据 \\
pypdf                     & 6.10.0  & RAG 知识库 PDF 解析 \\
python-docx               & 1.2.0   & RAG 知识库 DOCX 解析 \\
\bottomrule
\end{longtable}

\begin{warnbox}[体积优化]
admin-backend Dockerfile 显式从 \texttt{download.pytorch.org/whl/cpu} 安装 torch，
避免 pip 默认拉取约 2 GB 的 CUDA wheel，使最终镜像减小约 1.5 GB。
\end{warnbox}

\subsection{admin-frontend 镜像（多阶段：node:20-alpine $\rightarrow$ nginx:1.27-alpine）}

\textbf{Dockerfile 位置：}\texttt{admin/frontend/Dockerfile}

\begin{longtable}{L{4.0cm} L{2.0cm} L{7.0cm}}
\toprule
\textbf{依赖} & \textbf{版本} & \textbf{说明} \\
\midrule
\endhead
vue                        & \^{}3.5.32   & 核心框架 \\
vue-router                 & \^{}4.6.4    & 路由 \\
pinia                      & \^{}3.0.2    & 状态管理 \\
axios                      & \^{}1.15.0   & HTTP 客户端 \\
element-plus               & \^{}2.9.7    & 暗黑主题组件库 \\
@element-plus/icons-vue    & \^{}2.3.1    & 图标集 \\
echarts                    & \^{}6.0.0    & 图表库 \\
vue-echarts                & \^{}8.0.1    & ECharts Vue 封装 \\
three                      & \^{}0.183.2  & 3D 工厂可视化 \\
dayjs                      & \^{}1.11.13  & 时间格式化 \\
vite                       & \^{}5.4.19   & 构建工具（仅 build 阶段） \\
typescript                 & \~{}6.0.2    & 类型系统（仅 build 阶段） \\
vue-tsc                    & \^{}3.2.6    & Vue 类型检查（仅 build 阶段） \\
\bottomrule
\end{longtable}

% ========== 5 关键配置项 ==========
\section{关键配置项}

\subsection{\texttt{.env} 环境变量}

\begin{longtable}{L{4.8cm} L{3.5cm} L{5.5cm}}
\toprule
\textbf{变量} & \textbf{默认值} & \textbf{说明} \\
\midrule
\endhead
\texttt{ADMIN\_JWT\_SECRET}            & please-change-me-...      & JWT 签名密钥，生产必须更换 \\
\texttt{ADMIN\_USERNAME}               & admin                     & 管理员登录名 \\
\texttt{ADMIN\_PASSWORD}               & admin123456               & 管理员登录密码 \\
\texttt{ADMIN\_TOKEN\_EXPIRE\_MINUTES} & 120                       & Token 有效期（分钟） \\
\texttt{PG\_HOST}                      & 98.142.241.155            & 默认指向开发库，生产覆盖 \\
\texttt{PG\_PORT}                      & 5432                      & PostgreSQL 端口 \\
\texttt{PG\_DB}                        & aqimonitor                & 库名 \\
\texttt{PG\_USER}                      & team                      & 用户名 \\
\texttt{PG\_PASSWORD}                  & fwwb1234                  & 密码 \\
\texttt{OPENAI\_API\_KEY}              & （未设置）                 & RAG 大模型 Key（可选） \\
\texttt{OPENAI\_BASE\_URL}             & （未设置）                 & RAG 大模型地址（可选） \\
\bottomrule
\end{longtable}

\subsection{compose 内部固定项（不通过 .env 覆盖）}

\begin{longtable}{L{4.8cm} L{4.5cm} L{4.0cm}}
\toprule
\textbf{变量} & \textbf{值} & \textbf{说明} \\
\midrule
\endhead
\texttt{VOCS\_BASE\_URL}     & http://vocs-server:8001 & 容器内服务发现 \\
\texttt{ENSEMBLE\_BASE\_URL} & http://ensemble:8000    & 容器内服务发现 \\
\texttt{TZ}                  & Asia/Shanghai           & 时区 \\
\texttt{MODEL\_PATH}         & models/vocs\_seq2seq\_v2\_best.pth & ML 模型路径 \\
\texttt{SCALER\_PATH}        & models/vocs\_scalers\_v2.pkl       & 归一化器 \\
\bottomrule
\end{longtable}

\subsection{数据持久化}

\begin{itemize}[leftmargin=1.4em,itemsep=2pt]
  \item \texttt{./vocs\_realtime\_data} $\rightarrow$ \texttt{wg-vocs-server:/app/vocs\_realtime\_data}（CSV 回放数据）
  \item \texttt{./models} $\rightarrow$ \texttt{wg-vocs-server:/app/models}（模型文件，支持热更新）
  \item PostgreSQL 数据由\textbf{外部数据库}承载，不占用本机卷
\end{itemize}

% ========== 6 常见问题排查 ==========
\section{常见问题排查}

\subsection{容器启动失败 / 端口占用}

\begin{warnbox}[现象]
\texttt{docker compose up} 报错：\texttt{Bind for 0.0.0.0:3001 failed: port is already allocated}
\end{warnbox}

\textbf{排查与处理}：
\begin{enumerate}[leftmargin=1.6em,itemsep=2pt]
  \item Linux：\texttt{ss -ltnp | grep -E '3001|8000|8001|8003'} 找到占用进程；
  \item Windows：\texttt{netstat -ano | findstr :3001} 然后 \texttt{taskkill /PID <pid> /F}；
  \item 或修改 \texttt{docker-compose.prod.yml} 中端口映射，例如 \texttt{"3002:80"}。
\end{enumerate}

\subsection{前端白屏 / 控制台 \texttt{/api/* 502}}

\textbf{原因}：\texttt{admin-backend} 未启动或健康检查失败。

\begin{lstlisting}[language=bash]
docker compose -f docker-compose.prod.yml ps
docker compose -f docker-compose.prod.yml logs admin-backend --tail 100
# 常见根因：
#   - PG_PASSWORD 错误 → 看到 psycopg2.OperationalError
#   - VOCS_BASE_URL 配错 → 看到 httpx.ConnectError
\end{lstlisting}

\subsection{SSE 实时数据不刷新}

\textbf{排查清单}：
\begin{itemize}[leftmargin=1.4em,itemsep=2pt]
  \item \texttt{docker logs wg-vocs-server} 确认有数据接入；
  \item 浏览器 Network 面板查看 \texttt{/vocs/events} 是否处于 \texttt{pending} 状态；
  \item \texttt{nginx.conf} 中 \texttt{/vocs/} 与 \texttt{/api/} 已默认 \texttt{proxy\_buffering off}，
        如自定义反代需保留该配置；
  \item \texttt{proxy\_read\_timeout} 必须 $\geq$ 3600s，否则 60 秒后被动断开。
\end{itemize}

\subsection{ensemble 容器 OOM 被杀}

\begin{dangerbox}[现象]
\texttt{wg-ensemble} 状态变为 \texttt{Exited (137)}，\texttt{docker inspect} 中
\texttt{OOMKilled: true}。
\end{dangerbox}

\textbf{处理}：调大 compose 中 \texttt{deploy.resources.limits.memory}，或为宿主机加内存。
默认上限 4G 可满足常规推理；如有大批量并发请求建议提至 8G。

\subsection{PostgreSQL 连接失败}

\textbf{典型报错}：\texttt{psycopg2.OperationalError: could not connect to server}

\textbf{排查}：
\begin{enumerate}[leftmargin=1.6em,itemsep=2pt]
  \item 检查 \texttt{.env} 中 \texttt{PG\_HOST/PG\_USER/PG\_PASSWORD} 是否正确；
  \item 云数据库需将服务器出口 IP 加入白名单；
  \item 应用日志出现 \texttt{[DB] WARNING: init\_pool failed}——
        RAG 持久化功能将自动降级，但\textbf{不阻塞启动}，其他功能仍可用。
\end{enumerate}

\subsection{告警与持续监控面板"无数据"}

\begin{itemize}[leftmargin=1.4em,itemsep=2pt]
  \item 前端已内置 mock 演示数据（\texttt{stores/alerts.ts}、\texttt{stores/dashboard.ts}），
        即使后端断开也会回退显示；若 mock 也不出现，多为浏览器缓存了旧产物——按 Ctrl+F5 强刷。
  \item 确认 vocs-server 有数据：\texttt{curl http://localhost:8001/alerts?limit=5}。
\end{itemize}

\subsection{3D 场景卡顿 / WebGL 不可用}

\begin{itemize}[leftmargin=1.4em,itemsep=2pt]
  \item Chrome 地址栏输入 \texttt{about:gpu} 确认 WebGL2 处于 \texttt{Hardware accelerated}；
  \item 远程桌面 (RDP) 下 WebGL 常 fallback 到软件渲染——演示请用本机浏览器；
  \item 集成显卡建议浏览器窗口 $\leq$ 1920$\times$1080。
\end{itemize}

\subsection{构建慢 / 拉镜像超时}

\begin{lstlisting}[language=bash]
# 1) 配置 Docker 镜像加速器
sudo tee /etc/docker/daemon.json <<EOF
{
  "registry-mirrors": [
    "https://docker.mirrors.ustc.edu.cn",
    "https://hub-mirror.c.163.com"
  ]
}
EOF
sudo systemctl restart docker

# 2) admin-frontend Dockerfile 已默认 npm 国内源 (npmmirror.com)
# 3) Python 镜像可在 Dockerfile 取消下面这行的注释：
#    RUN pip config set global.index-url https://pypi.tuna.tsinghua.edu.cn/simple
\end{lstlisting}

\subsection{升级 / 回滚}

\textbf{升级（已推送到镜像仓库时）}：

\begin{lstlisting}[language=bash]
docker compose -f docker-compose.prod.yml pull
docker compose -f docker-compose.prod.yml up -d
\end{lstlisting}

\textbf{升级（源码方式）}：

\begin{lstlisting}[language=bash]
git pull
docker compose -f docker-compose.prod.yml build
docker compose -f docker-compose.prod.yml up -d
\end{lstlisting}

\textbf{回滚}：保留旧镜像 tag，例如 \texttt{:1.0.0} 与 \texttt{:1.1.0} 并存，
回滚时只需 \texttt{docker tag waste-gas/xxx:1.0.0 waste-gas/xxx:latest \&\& docker compose up -d}。

% ========== 7 附录 ==========
\section{附录}

\subsection{交付物清单}

\begin{longtable}{L{6.0cm} L{8.0cm}}
\toprule
\textbf{文件 / 目录} & \textbf{用途} \\
\midrule
\endhead
\texttt{docker-compose.prod.yml}              & 一键编排（4 个服务） \\
\texttt{.env.example}                         & 环境变量样板 \\
\texttt{.dockerignore}                        & 优化构建上下文 \\
\texttt{Dockerfile}                           & vocs-server 镜像 \\
\texttt{vocs\_server.py}                      & ML 推理主入口 \\
\texttt{models/}                              & 预训练模型（运行时挂载） \\
\texttt{vocs\_realtime\_data/}                & CSV 回放数据（运行时挂载） \\
\texttt{admin/backend/Dockerfile}             & admin-backend 镜像 \\
\texttt{admin/frontend/Dockerfile}            & admin-frontend 镜像（多阶段） \\
\texttt{admin/frontend/nginx.conf}            & 反向代理 + SSE 配置 \\
\texttt{models/new\_VOC/ensemble\_docker/}    & ensemble 镜像构建上下文 \\
\bottomrule
\end{longtable}

\subsection{健康检查命令速查}

\begin{lstlisting}[language=bash]
# 全量容器状态
docker compose -f docker-compose.prod.yml ps

# 各服务健康端点
curl -s http://localhost:8000/health        | python -m json.tool
curl -s http://localhost:8001/status        | python -m json.tool
curl -s http://localhost:8003/api/health
curl -s http://localhost:3001/healthz

# 容器内部网络互通验证
docker exec wg-admin-backend curl -s http://vocs-server:8001/status
docker exec wg-admin-backend curl -s http://ensemble:8000/health
\end{lstlisting}

\subsection{日志位置}

\begin{itemize}[leftmargin=1.4em,itemsep=2pt]
  \item Docker 全栈日志：\texttt{docker compose -f docker-compose.prod.yml logs -f}
  \item 单服务日志：\texttt{docker logs -f wg-admin-backend}
  \item Nginx 访问日志（前端容器内）：\texttt{/var/log/nginx/access.log}
\end{itemize}

\subsection{与 v1.0 手册的差异}

\begin{longtable}{L{4.5cm} L{4.5cm} L{4.5cm}}
\toprule
\textbf{条目} & \textbf{v1.0（源码部署）} & \textbf{v2.0（Docker 部署）} \\
\midrule
\endhead
部署步骤数      & $\geq$ 10 步         & \textbf{4 步} \\
依赖安装        & 手动装 Python/Node/Nginx 等 & 仅需 Docker \\
环境冲突风险    & 高（torch 版本、libpq 等） & 已隔离 \\
首次部署耗时    & 30-60 分钟            & 10-20 分钟（多数为镜像构建） \\
离线交付        & 需打包源码 + 文档教安装  & \texttt{docker save/load} 一键 \\
回滚            & 手动逐项恢复          & 切换镜像 tag 一键 \\
\bottomrule
\end{longtable}

\vspace{1em}
\begin{infobox}[技术支持]
如排查未果，请联系项目组并提供：
\begin{enumerate}[leftmargin=1.4em,itemsep=1pt]
  \item \texttt{docker compose -f docker-compose.prod.yml ps} 截图
  \item 问题服务最近 200 行日志（\texttt{docker logs --tail 200 <container>}）
  \item 浏览器控制台 Network 面板截图（如为前端问题）
  \item Docker / Docker Compose 版本（\texttt{docker version}、\texttt{docker compose version}）
\end{enumerate}
\end{infobox}

\vspace{2em}
\begin{center}
\color{gray}\small --- 文档结束 ---\\[2pt]
智洁园区项目组 \textperiodcentered{} 2026
\end{center}

\end{document}
